\documentclass[11pt]{amsart}
\usepackage{amsmath,amssymb,graphicx,hyperref,longtable}
\newtheorem{theorem}{Theorem}
\newtheorem{conjecture}[theorem]{Conjecture}
\newtheorem{proposition}[theorem]{Proposition}
\theoremstyle{remark}
\newtheorem{remark}[theorem]{Remark}

\title[The smoothed staple]{The smoothed staple: shorter certified escape
paths for Bellman's lost-in-a-forest problem in isosceles triangles}
\author{Alessio Doria}
\date{July 2026 --- working draft}

\thanks{\textbf{Status.} This is an informal working draft, circulated for
verification only. It is not submitted, and not intended for submission in
its present form. Its purpose is to state precisely what the accompanying
machine-checkable certificates verify, and on what imported result they
depend, so that the argument can be checked by others. The reasoning behind
the certificates --- in particular the variant of the verification scheme
described in Section~2 --- has not been widely reviewed.
Comments and corrections are very welcome.}

\begin{document}
\begin{abstract}
For the isosceles triangle $T$ with unit legs and base angle $\beta$,
Temerev (2026) proved the certified upper bound
$\ell(T)\le 1.2849615334\ldots$ at the golden gnomon
($\beta=36^\circ$) via a three-segment ``staple'' path, together with an
interval theorem for $\beta\in[31.08^\circ,42.16^\circ]$. We show that, at the
golden gnomon, the three-segment staple is not optimal among
convex-position polygonal paths satisfying the same escape criterion. Using the same exact reduction
(escape of a path follows from a single polygon-contains-disk inclusion for
a weighted Minkowski sum $W$ of rotated copies of the path's convex hull),
we exhibit an explicit $20$-vertex path with rational coordinates proving
\[
  \ell(T)\;\le\;1.2826879686\ldots\;<\;1.2849615334\ldots
\]
at the golden gnomon, with every accepted sign determination carried out in
exact arithmetic over $\mathbb{Q}(\sqrt5,\sqrt{10-2\sqrt5})$. Numerical
optimization suggests that the underlying continuum object is a symmetric
\emph{smoothed staple} whose shoulders are circular arcs, and a bang-bang
argument for the
associated infinite-dimensional linear program in support functions leads us
to conjecture that their curvature radius is exactly $\sin\beta$.
Conservative per-angle computations indicate the family beats the scaled
Zalgaller caliper down to $\beta\approx30.76^\circ$ (staple:
$31.05^\circ$). Optimality is not claimed.
\end{abstract}
\maketitle

\section{Introduction}
Bellman's lost-in-a-forest problem asks for the shortest path guaranteed to
reach the boundary of a convex region of known shape from an unknown
starting point and heading \cite{FW04}. For isosceles triangles with base
angles below $45^\circ$ the problem is open; Ward \cite{Ward08} compared the
classical candidates numerically, and Temerev \cite{Tem26} proved the first
certified improvement on an interval of base angles via a piecewise-linear
three-segment family of \emph{staples}, including the bound
$\ell(T)\le 1.2849615334\ldots$ at the golden gnomon
($\beta=36^\circ$, unit legs).

This note improves those bounds. The key observation is structural: the
escape criterion used in \cite{Tem26} depends only on the convex hull of the
path, so the staple is a single point in the much larger space of paths with
convex-position vertices. Optimizing freely in that space, under the
\emph{same exact criterion}, produces strictly shorter escaping paths, and
the limiting shape is not piecewise linear.

\section{Preliminaries: the exact escape reduction}
Let $T=T(\beta)$ have unit legs, base angle $\beta$, outward edge normals
$n_0=(0,-1)$, $n_{1,2}=(\pm\sin\beta,\cos\beta)$ with weights (edge lengths)
$w_0=2\cos\beta$, $w_1=w_2=1$, and let $\rho=2\,\mathrm{Area}(T)=\sin2\beta$.
For a convex body $H$ put
\[
  W(H)\;=\;\sum_{i=0}^{2} w_i\,R_{n_i}H,
\]
the Minkowski sum of rotated copies of $H$, where $R_n$ is the rotation
taking $(1,0)$ to $n$'s frame as in \cite{Tem26}. Since
$\sum_i w_i R_{n_i}=0$, the construction is translation invariant, and since
plane rotations commute, the criterion below is rotation invariant.

\begin{proposition}[imported from \cite{Tem26}, via Lutwak \cite{Lut98}]
\label{prop:reduction}
If $\operatorname{disk}(\rho)\subseteq W(H)$, where $H$ is the convex hull
of a path $P$, then $P$ escapes $T(\beta)$.
\end{proposition}

All certificates below verify the inclusion of
Proposition~\ref{prop:reduction} for an explicit polygon, by exhibiting a
convex, positively oriented cycle of points of $W$ whose edge lines all lie
at distance $\ge\rho$ from the origin. Because the cycle's vertices belong
to $W$ and $W$ is convex, the cycle's hull is contained in $W$; hence the
disk inclusion follows from the finitely many edge conditions
$\mathrm{cross}(p,q)^2-\rho^2\,|q-p|^2>0$, each an exact sign determination
in $\mathbb{Q}(\sqrt5,\sqrt{10-2\sqrt5})$ at $\beta=36^\circ$.

\section{Main theorem}
\begin{theorem}\label{thm:main}
Let $T$ be the golden gnomon (apex $108^\circ$, base angles $36^\circ$,
unit legs). The $20$-vertex polygonal path with the rational vertices of
Appendix~\ref{tab:verts} (traversed in order) escapes $T$, and consequently
\[
  \ell(T)\;\le\;1.2826879686\ldots\;<\;1.2849615334\ldots
\]
\end{theorem}

\begin{proof}[Proof (machine-checked)]
The escape inclusion is verified by
\texttt{verification/verify\_certificate20.py} in the accompanying
repository: the proposed $57$-vertex cycle of Minkowski points satisfies,
for every edge, positive convexity cross product, positive orientation
cross product, and $D=\mathrm{cross}^2-\rho^2|e|^2>0$, all by exact
symbolic sign determinations (\texttt{sympy}); the minimum edge distance is
$0.95105745\ldots>\rho=\sin72^\circ=0.95105651\ldots$. The path length, a
sum of $19$ square roots of rationals, is bounded above by exact rational
majorants of each square root (integer arithmetic), and the resulting
rational is compared exactly against $12826879687/10^{10}$. The certificate
cycle is stored explicitly with the data; floating point appears in no
accepted check.
\end{proof}

\begin{remark}
A shorter, fully self-contained $8$-vertex certificate proving
$\ell(T)\le1.2827871$ is included for independent auditing
(\texttt{smoothed\_staple\_certificate.py}).
\end{remark}

\section{The continuum object and the $\sin\beta$ curvature law}
Free polyline optimization under the exact criterion yields the strictly
decreasing values $1.2849609$ ($n=4$, the staple), $1.2829755$ ($n=6$),
$1.2827742$ ($n=8$), $1.2827491$ ($n=10$). A piecewise line/arc model with
free radii, solved by constraint generation on the support-function form of
the inclusion, converges (dense margin $-6\cdot10^{-15}$) to the symmetric
path
\[
  \underbrace{0.291027}_{\text{straight}}
  \;\to\;\underbrace{r=0.587785,\ \varphi=9.5232^\circ}_{\text{arc}}
  \;\to\;\underbrace{0.071683}_{\text{straight}}
  \;\to\;\underbrace{76.3054^\circ}_{\text{corner}}
  \;\to\;\underbrace{0.361665}_{\text{base}}\;\to\;(\text{mirror}),
\]
of total length $L_\infty=1.2826726243$. The arc radius agrees with
$\sin36^\circ=0.5877853$ to the optimizer's precision, with the radius left
free; a second arc allowed on each shoulder degenerates.

\emph{Relation to prior work.} The line--arc geometry is closely related to
Gibbs's ``Tunnel'' solutions, proposed on Monte Carlo evidence for base
angles roughly $30.8^\circ$--$38.1^\circ$ \cite{Gib16a} --- a range
containing $36^\circ$ --- and coupled curvature conditions for extremal
escape paths appear in his analysis for convex polygons \cite{Gib16b}.
Those proposals were not accompanied by certified bounds. The contribution
of the present note is therefore \emph{not} the qualitative smoothed
geometry, which should be credited to Gibbs, but the exact rational
polygonal certificate showing that a smoothed/Tunnel-type path rigorously
improves the previously certified staple bound at the golden gnomon, and
the specific quantitative prediction $r=\sin\beta$ below.

The value $\sin\beta$ is predicted by the following argument. The
constraint of Proposition~\ref{prop:reduction} reads
$h_W(u)=\sum_i w_i\,h_H(R_{n_i}^{T}u)\ge\rho$ for all directions $u$ and is
\emph{linear} in the support function $h_H$; path length is likewise a
linear functional of $h_H$ over the traversed normal window. The problem is
therefore an infinite-dimensional linear program. On a window where the
constraint is active, $h_W\equiv\rho$ gives $h_W+h_W''=\rho$, i.e.
\[
  \sum_i w_i\,\bigl(h_H+h_H''\bigr)(\theta-\alpha_i)\;=\;\rho ,
\]
so the curvature radii of $\partial H$ at the three rotated angles combine,
with weights $w_i$, to the constant $\rho$. A bang-bang optimum concentrates
all curvature on the copy with the largest weight $w_0=2\cos\beta$, giving
curvature radius
\[
  \frac{\rho}{2\cos\beta}\;=\;\frac{\sin2\beta}{2\cos\beta}\;=\;\sin\beta .
\]

\begin{conjecture}\label{conj:curvature}
For $\beta$ in the staple regime, the optimal escape path of $T(\beta)$ is a
concatenation of line segments, corners, and circular arcs of curvature
radius exactly $\sin\beta$, symmetric about the axis of $T$; at
$\beta=36^\circ$ its length is $1.2826726\ldots$
\end{conjecture}

\section{Per-angle behaviour and the caliper crossing}
Conservative per-angle values (grid-optimized $8$-vertex polylines followed
by a conservative scaling correction evaluated in floating point against
the hull-edge criterion; unlike the certificates of Section~3 these
corrections are numerical, not symbolic) give a uniform gain of $\approx+0.0026$ over
the per-angle optimal staple for $\beta\in[30^\circ,33^\circ]$, moving the
crossing with the scaled Zalgaller caliper from $\beta\approx31.05^\circ$
(staple) to
\[
  \beta\;\approx\;30.76^\circ\qquad(\text{conservative; the continuum
  crossing is lower still}).
\]
The family also beats the low-regime square path at every sampled angle down
to $29^\circ$; square comparisons are made only below Ward's regime switch
$\beta_0\approx39.13^\circ$. Near the upper end the smoothing gain genuinely collapses
($+1.6\cdot10^{-4}$ at $42.0^\circ$, $+10^{-5}$ at $42.5^\circ$; $n=10$
polylines with exact-margin correction, the same recipe that captures
$\approx90\%$ of the $2.3\cdot10^{-3}$ gain at $36^\circ$): the arcs
switch off as $\beta$ approaches the zee crossing, which therefore stays at
$\beta\approx42.284^\circ$, essentially the staple's $42.28^\circ$ and
short of the conjectured $42.3^\circ$.

\section{Open problems}
(i) An interval theorem for the smoothed family in $u=\tan(\beta/2)$,
replacing arcs by circumscribed tangent polygons and reusing the Sturm
certificate machinery of \cite{Tem26}; (ii) the exact KKT system at the
gnomon with $r=\sin\beta$ imposed (five unknowns), which may yield a closed
form for $1.2826726\ldots$; (iii) a proof of
Conjecture~\ref{conj:curvature}; (iv) matching lower bounds.

\appendix
\section{Certificate vertices}\label{tab:verts}
The $20$ vertices (path order) of the Theorem~\ref{thm:main} witness;
denominators $10^8$ or divisors thereof.
\begin{longtable}{lll}
\hline
vertex & $x$ & $y$\\
\hline
$v_{0}$ & $-1256697/12500000$ & $9697297/25000000$ \\
$v_{1}$ & $-18524027/100000000$ & $11076737/100000000$ \\
$v_{2}$ & $-18800781/100000000$ & $10171311/100000000$ \\
$v_{3}$ & $-600917/3125000$ & $4309011/50000000$ \\
$v_{4}$ & $-4903679/25000000$ & $1763883/25000000$ \\
$v_{5}$ & $-99787/500000$ & $5482169/100000000$ \\
$v_{6}$ & $-20257379/100000000$ & $155879/4000000$ \\
$v_{7}$ & $-20513533/100000000$ & $2305271/100000000$ \\
$v_{8}$ & $-10785997/50000000$ & $-5656137/100000000$ \\
$v_{9}$ & $-5394019/25000000$ & $-568721/10000000$ \\
$v_{10}$ & $12159143/100000000$ & $-9379621/50000000$ \\
$v_{11}$ & $193587/1562500$ & $-18428471/100000000$ \\
$v_{12}$ & $8348493/50000000$ & $-12229211/100000000$ \\
$v_{13}$ & $17464737/100000000$ & $-11066231/100000000$ \\
$v_{14}$ & $9100817/50000000$ & $-4945211/50000000$ \\
$v_{15}$ & $9453909/50000000$ & $-8702301/100000000$ \\
$v_{16}$ & $19584521/100000000$ & $-3750221/50000000$ \\
$v_{17}$ & $5057973/25000000$ & $-6284817/100000000$ \\
$v_{18}$ & $20847673/100000000$ & $-2530147/50000000$ \\
$v_{19}$ & $33614857/100000000$ & $10934787/50000000$ \\
\hline
\end{longtable}

\begin{thebibliography}{9}
\bibitem{FW04} S.~R. Finch, J.~E. Wetzel, \emph{Lost in a forest}, Amer.
Math. Monthly 111 (2004), 645--654.
\bibitem{Gib16a} P.~E. Gibbs, \emph{Lost in an Isosceles Triangle},
working paper (July 2016).
\bibitem{Gib16b} P.~E. Gibbs, \emph{Bellman's Escape Problem for Convex
Polygons}, working paper (June 2016).
\bibitem{Lut98} E.~Lutwak, \emph{Containment and circumscribing simplices},
Discrete Comput. Geom. 19 (1998), 229--235.
\bibitem{Tem26} A.~Temerev, \emph{An improved upper bound for Bellman's
lost-in-a-forest problem in isosceles triangles} (2026),
\url{https://github.com/atemerev/bellman-staple}.
\bibitem{Ward08} J.~W. Ward, \emph{Exploring the Bellman forest problem},
manuscript (2008).
\end{thebibliography}
\end{document}
